% status: FULL
% Authoritative LaTeX; edit this file, not the historical Markdown.
\chapter{Matrix Multiplication: The Engine Room}\label{chapter-6-matrix-multiplication-the-engine-room}\label{app:matmul}
\Appref{tensors} gave tensor values addresses. Shape said which logical indices
were legal, strides mapped those indices to storage, and ownership said
how long the bytes remained alive. Yet an inference engine exists to
\emph{compute} with those values. Its most familiar numerical operation
is also one of its most consequential systems problems: matrix
multiplication.

A mathematical line such as \(C=AB\) looks atomic. A processor does not
receive that line. It receives loads, multiplications, additions,
stores, branches, and addresses. The same mathematical products can be
visited in different orders. Those orders can move very different
amounts of data through a memory hierarchy, even when their operation
counts are identical.

This chapter builds the book's first optimization cycle:

A model operation declares a mathematical result; a checked kernel maps its indices to storage; a measured candidate may replace that reference only after equivalence checks.

The cycle matters more than the particular optimization. ENGINE-2 will
contain an obviously inspectable scalar reference path and a
cache-blocked scalar CPU path. It will keep both. The reference defines
behavior for general valid views; the blocked path accepts a narrower
layout contract and must earn its place with correctness gates and
measurements.

\begin{quote}
\textbf{FIRST PRINCIPLE} Matrix multiplication specifies which values
contribute to an answer. A kernel specifies the order in which
arithmetic and byte movement realize that answer.
\end{quote}

\EngineFigure{ch06-engine}{From equation through addresses, reference loops, reuse, and hardware to model output.}{fig:ch06-engine}

Read this chapter in three passes: first determine \textbf{which values}
contribute; then trace \textbf{which addresses} the implementation
visits; finally ask \textbf{what the measurements justify}. A fast wrong
answer fails the first pass. A plausible cache story without timing
evidence fails the third. The production comparison at the end
translates these same questions into Hermon and GGML, without adding a
production backend to the teaching engine.

\section{The most important loop in
AI}\label{the-most-important-loop-in-ai}

Language-model inference repeatedly applies learned linear
transformations. Later chapters will add normalization, position
handling, attention, and feed-forward networks. Under many of those
operators sit matrix-vector or matrix-matrix products. They project
activations from one feature space into another; they also dominate
large regions of weight storage and execution.

That importance does not make every product the same. During one-token
decode, a weight matrix may multiply one activation vector. During work
over several tokens or requests, the same weights may multiply several
columns or rows of activations. The equations are closely related, but
their opportunity to reuse each weight after loading it is different.
Shape is therefore part of the performance problem, not merely a type
check.

ENGINE-1 already contained a linear projection:

\begin{equation}
\mathbf{z}=\mathbf{W}\mathbf{h}+\mathbf{b},
\end{equation}

where \(\mathbf{W}\in\mathbb{R}^{V_{\mathrm{vocab}}\times D}\),
\(\mathbf{h}\in\mathbb{R}^{D}\),
\(\mathbf{b}\in\mathbb{R}^{V_{\mathrm{vocab}}}\), and
\(\mathbf{z}\in\mathbb{R}^{V_{\mathrm{vocab}}}\). The old implementation
spelled the matrix-vector product as loops inside the model. That was
appropriate when the loop first made logits real. It is now duplication.
\Appref{matmul} extracts the numerical operation into a kernel layer while
leaving embedding lookup and bias addition as explicit model operations.

\section{From multiply to
multiply-accumulate}\label{from-multiply-to-multiply-accumulate}

Begin with one multiplication. Given two \texttt{f32} values \(a\) and
\(b\), their product contributes \(ab\) to some destination. A linear
algebra kernel rarely needs that product alone. It adds many products
into an accumulator:

\begin{equation}
s \leftarrow s + ab.
\end{equation}

This is a multiply-accumulate step. Hardware may later implement the
expression with a fused multiply-add instruction, which rounds once
rather than rounding the product and sum separately. ENGINE-2 does not
select intrinsics or promise fused behavior. Its source uses
\texttt{f32} operands and an \texttt{f32} accumulator. The independent
oracle models separate \texttt{f32} rounding, while cross-path tests use
an explicit tolerance so compiler and architecture choices cannot
masquerade as semantic failure.

Follow one contribution in the canonical
\href{https://github.com/hermonai/inside-the-llm-engine/blob/astra-visual-rewrite/diagrams/linear/follow-the-flop.txt}{FLOP
diagram}:

\begin{equation}C_{ij}=\sum_k \underbrace{A_{ik}}_{\text{row }i}\underbrace{B_{kj}}_{\text{column }j}.\end{equation}

The floating-point representation makes the update bounded in precision,
not exact over real numbers; overflow can produce infinity. If large
positive and negative contributions nearly cancel, rounding error can be
visible. If an optimized kernel changes the reduction order, its final
low bits may change as well.

\section{Dot product}\label{dot-product}
\index{dot product}

For two vectors \(\mathbf{a},\mathbf{b}\in\mathbb{R}^{K}\), the dot
product is

\begin{equation}
\mathbf{a}\cdot\mathbf{b}
=\sum_{k=0}^{K-1}a_kb_k.
\end{equation}

For the chapter's canonical example, take \(\mathbf{a}=[1,-2,3,0]\) and
\(\mathbf{x}=[2,1,-1,0.5]\):

\begin{equation}
\mathbf{a}\cdot\mathbf{x}
=1(2)+(-2)(1)+3(-1)+0(0.5)
=2-2-3+0=-3.
\end{equation}

\EngineFigure{ch06-dot}{Four matched-index products reduce to one scalar, minus three.}{fig:ch06-dot}

The two operands enter each multiplication independently. The product at
index 2 is \(3(-1)\), never \(3(0.5)\): matching the reduction index is
the operation's identity. A visually pleasing crossing arrow that pairs
the wrong indices would teach a different computation.

The
\href{https://github.com/hermonai/inside-the-llm-engine/blob/astra-visual-rewrite/diagrams/linear/dot-product-multiply-accumulate.txt}{multiply-accumulate
diagram} shows the reduction visually. The shape contract is strict:
both operands have tensor rank 1 and the same length \(K\). ENGINE-2
does not silently truncate to the shorter length or broadcast one
element.

When \(K=0\), the loop has no terms. The sum is the additive identity,
zero. This is not a special numerical trick; it is the standard
empty-sum convention and makes zero-dimensional matrix products
composable.

\section{Build it: scalar dot}\label{build-it-scalar-dot}

The reference implementation is deliberately direct:

\begin{verbatim}
let mut sum = 0.0_f32;
for index in 0..left_len {
    sum += *left.get(&[index])? * *right.get(&[index])?;
}
\end{verbatim}

Validation happens before the loop. \texttt{TensorView::get} then
follows each operand's logical indexing contract. A stride-2 vector and
a zero-stride read-only broadcast view are therefore valid. The kernel
is not assuming that logical neighbors occupy adjacent storage.

\begin{quote}
\textbf{BUILD IT} Complete
\href{https://github.com/hermonai/inside-the-llm-engine/blob/astra-visual-rewrite/labs/lab-22-dot-product.md}{Lab
22}. Expand the arithmetic, run \texttt{dot\_reference}, then demand
typed failures for bad rank and length.
\end{quote}

This function is not meant to beat a tuned vector library. It gives
later kernels a legible semantic anchor: one reduction, ordered from
index 0 through \(K-1\), into \texttt{f32}.

\section{Matrix times vector}\label{matrix-times-vector}

Let \(\mathbf{A}\in\mathbb{R}^{M\times K}\) and
\(\mathbf{x}\in\mathbb{R}^{K}\). Their matrix-vector product
\(\mathbf{y}=\mathbf{A}\mathbf{x}\in\mathbb{R}^{M}\) has one dot product
per matrix row:

\begin{equation}
y_i=\sum_{k=0}^{K-1}A_{ik}x_k,
\qquad 0\le i<M.
\end{equation}

The shared \(K\) dimension is contracted. The \(M\) row dimension
survives into the output. The complete
\href{https://github.com/hermonai/inside-the-llm-engine/blob/astra-visual-rewrite/diagrams/linear/gemv-shape-contract.txt}{GEMV
shape diagram} makes that flow explicit.

Use one rectangular fixture throughout the numerical plates. We call the
weight matrix \(W\) here and use \(A=W\) when discussing generic GEMM:

\begin{equation}
W=
\begin{bmatrix}
1&-2&3&0\\
0&1&-1&2\\
2&0&1&-1
\end{bmatrix},
\qquad
\mathbf{x}=
\begin{bmatrix}
2\\1\\-1\\0.5
\end{bmatrix},
\end{equation}

the first row produces

\begin{equation}
y_0=1(2)+(-2)(1)+3(-1)+0(0.5)=-3,
\end{equation}

while the remaining rows produce

\begin{equation}
\begin{aligned}
y_1&=0(2)+1(1)+(-1)(-1)+2(0.5)=3,\\
y_2&=2(2)+0(1)+1(-1)+(-1)(0.5)=2.5.
\end{aligned}
\end{equation}

Thus \(\mathbf{y}=[-3,3,2.5]\). The rectangular matrix and mixed signs
help expose orientation errors that square, all-positive fixtures can
conceal.

\EngineFigure{ch06-row}{Four increasing-k accumulator states end at y zero equals minus three.}{fig:ch06-row}

The partial sums are \(2,0,-3,-3\), starting from zero before the first
product. The final zero product still occupies a legitimate reduction
index. The
\href{https://github.com/hermonai/inside-the-llm-engine/blob/astra-visual-rewrite/figures/generated/ch06-row.html}{step-through
row animation} highlights each update; this static plate contains every
state and loses no information in print. Nothing starts moving
automatically.

\EngineFigure{ch06-gemv}{The same x contributes independently to three row dot products and three output values.}{fig:ch06-gemv}

\section{Build it: GEMV}\label{build-it-gemv}

\texttt{gemv\_reference} accepts a rank-2 matrix view and a rank-1
vector view. It validates \(A.shape[1]=x.shape[0]\), allocates a new
canonical \texttt{{[}M{]}} \texttt{OwnedTensor}, and fills it with row
dot products. Inputs remain immutable and borrowed; output storage
belongs to the caller.

The ownership choice is important. Returning a view into a temporary
work buffer would entangle result lifetime with kernel internals.
Mutating an input as an output would introduce aliasing questions. A
fresh owner makes the v1 contract simple:

The two inputs are immutable views; the returned $[M]$ result owns a separate allocation. A result therefore survives the return from the kernel without borrowing its temporary workspace.

The numerical body is the existing kernel, not a diagram-only
implementation:

\begin{verbatim}
for (row, value) in output.iter_mut().enumerate() {
    let mut sum = 0.0_f32;
    for k in 0..inner {
        sum += *matrix.get2(row, k)? * *vector.get(&[k])?;
    }
    *value = sum;
}
\end{verbatim}

For shape \texttt{{[}M,0{]}\ ×\ {[}0{]}}, the result is an owned
\texttt{{[}M{]}} vector of zeros. For \texttt{{[}0,K{]}\ ×\ {[}K{]}}, it
is an empty owner with shape \texttt{{[}0{]}}. No sentinel or fallback
path is needed.

\begin{quote}
\textbf{BUILD IT} Complete
\href{https://github.com/hermonai/inside-the-llm-engine/blob/astra-visual-rewrite/labs/lab-23-gemv-by-hand.md}{Lab
23}, including the strided fixture. Physical padding must not affect the
logical result.
\end{quote}

\section{From a weight coordinate to a loaded
F32}\label{from-a-weight-coordinate-to-a-loaded-f32}

\Appref{tensors}'s address contract remains active inside every multiply. For a
valid rank-2 view with element strides \(s_0,s_1\) and base offset
\(o\),

\begin{equation}
\operatorname{offset}_W(i,k)=o+i s_0+k s_1.
\end{equation}

Our canonical \(W:[3,4]\) has strides \texttt{{[}4,1{]}} and base zero.
Therefore \(W_{1,2}\) lives at element \(1(4)+2=6\), whose displacement
from the allocation's start is \(4(6)=24\) bytes for F32. The loaded
value is \(-1\). An element index, a byte displacement, and the value at
that location are three different things.

\EngineFigure{ch06-memory}{Logical W maps to twelve F32 storage cells; index one comma two is offset six and byte twenty-four.}{fig:ch06-memory}

The reference calls checked view access; it does not replace that
mapping with \texttt{i*K+k} for arbitrary inputs. A transposed view
changes strides without moving values. A zero-stride immutable view may
reuse one storage cell at several logical positions. Both can be valid
reference operands. The blocked kernel later requires canonical
row-major slices so its simplified arithmetic is justified by an
explicit boundary check, not by an assumption hidden in a loop.

The canonical
\href{https://github.com/hermonai/inside-the-llm-engine/blob/astra-visual-rewrite/code/reference/fixtures/chapter06-visual.json}{fixture}
stores inputs only. The independent
\href{https://github.com/hermonai/inside-the-llm-engine/blob/astra-visual-rewrite/code/reference/python/chapter06_visual_oracle.py}{Python
visual oracle} derives outputs, addresses, contribution traces, and tile
ranges. The
\href{https://github.com/hermonai/inside-the-llm-engine/blob/astra-visual-rewrite/code/mini-engine/crates/engine0/examples/chapter06_visual_trace.rs}{Rust
visual example} uses the real tensor and linear APIs. A parity script
compares the results; neither a manually typed answer nor the artwork
alone is the specification.

\section{Matrix times matrix}\label{matrix-times-matrix}

Let \(\mathbf{A}\in\mathbb{R}^{M\times K}\) and
\(\mathbf{B}\in\mathbb{R}^{K\times N}\). Their matrix product
\(\mathbf{C}=\mathbf{A}\mathbf{B}\in\mathbb{R}^{M\times N}\) has
elements

\begin{equation}
C_{ij}=\sum_{k=0}^{K-1}A_{ik}B_{kj},
\qquad 0\le i<M,\quad 0\le j<N.
\end{equation}

Every output cell is the dot product of row \(i\) from \(A\) and column
\(j\) from \(B\). The inner dimensions must agree. There is no
requirement that \(M\), \(K\), and \(N\) equal one another.

The two matrices are independent inputs to one contraction. An arrow
from \(A\) into \(B\) would incorrectly suggest that the first produces
the second.

See the canonical
\href{https://github.com/hermonai/inside-the-llm-engine/blob/astra-visual-rewrite/diagrams/linear/gemm-shape-contract.txt}{GEMM
shape contract} and
\href{https://github.com/hermonai/inside-the-llm-engine/blob/astra-visual-rewrite/diagrams/linear/gemm-one-output-cell.txt}{one-output-cell
diagram}. The common acronym \textbf{GEMM} comes from the BLAS general
matrix-matrix family; \textbf{GEMV} names the corresponding
matrix-vector family. ENGINE-2 implements a much smaller contract than
full BLAS: no transpose flags, scaling parameters, mixed layouts, or
in-place accumulation into caller storage.

\section{Shape contracts are execution
contracts}\label{shape-contracts-are-execution-contracts}
\index{matrix multiplication!shape contract}

The symbols \(M\), \(K\), and \(N\) are not decorative labels. They
control loop bounds, allocation size, and address arithmetic. Each
dimension has a role:

\begin{itemize}
\tightlist
\item
  \(M\) counts independent rows of the left operand and output;
\item
  \(K\) counts contributions reduced into each output value;
\item
  \(N\) counts columns of the right operand and output.
\end{itemize}

The agreement rule is positional. For \(A:[M,K]\) and \(B:[K,N]\), the
last axis of \(A\) equals the first axis of \(B\). Equal total element
counts are irrelevant. Shapes \texttt{{[}2,6{]}} and \texttt{{[}3,4{]}}
both contain twelve values, but they cannot occupy the two sides of this
product because 6 does not equal 3. Reshaping one operand would change
the mathematical problem and must be an explicit caller action.

The dimension roles also explain orientation. ENGINE-1 stores projection
weights as \texttt{{[}vocabulary,\ hidden{]}}. Vocabulary candidates are
output rows, so \(M=V\) and \(K=D\). Multiplying by hidden
\texttt{{[}D{]}} naturally returns \texttt{{[}V{]}}. If the weights were
stored \texttt{{[}D,V{]}}, the call would require a transpose view or a
different contract. A kernel must never guess orientation from equal
numbers.

Shape validation precedes output allocation for two reasons. First, it
avoids spending memory on a computation that has no defined result.
Second, it keeps failure deterministic: an inner-dimension mismatch
remains that typed error rather than sometimes becoming allocation
failure first. Output element count \(MN\) is checked independently
because valid input views with empty K can still describe an
unrepresentably large prospective output.

Ranks are checked before individual axes are read.
\texttt{gemv\_reference} cannot interpret a rank-2 shape as a vector
merely because it contains one column;
\texttt{\seqsplit{matmul\_reference}} cannot flatten a higher-rank
tensor. Later engines will define batching explicitly. Treating extra
axes as implicit batches here would create a second, undocumented
operator.

This strictness makes error messages part of the kernel API.
\texttt{RankMismatch} names the operation and operand.
\texttt{\seqsplit{InnerDimensionMismatch}} reports both K values.
\texttt{\seqsplit{OutputShapeOverflow}} reports M and N. Callers can
diagnose which contract failed without reconstructing it from a generic
bounds error.

\section{One output cell by hand}\label{one-output-cell-by-hand}

Keep \(A=W\) from GEMV, and place two input vectors into the columns of
\(B\):

\begin{equation}
B=
\begin{bmatrix}
2&1\\
1&-1\\
-1&2\\
0.5&0
\end{bmatrix}.
\end{equation}

The selected output cell follows row 0 and column 1, matching the same
\(k\):

\begin{equation}
C_{01}=1(1)+(-2)(-1)+3(2)+0(0)=9.
\end{equation}

Therefore

\begin{equation}
C=
\begin{bmatrix}
-3&9\\
3&-3\\
2.5&4
\end{bmatrix}.
\end{equation}

\EngineFigure{ch06-gemm}{Selected row and column independently feed the dot product for C zero comma one, with column and outer-product interpretations.}{fig:ch06-gemm}

The same contraction has two other useful interpretations. First, each
column of \(B\) is a GEMV input:

\begin{equation}
C_{:,j}=A B_{:,j},\qquad B_{:,0}=\mathbf{x},\qquad C_{:,0}=\mathbf{y}.
\end{equation}

Second, one reduction index contributes a whole rank-one matrix:

\begin{equation}
C=\sum_{k=0}^{K-1} A_{:,k}B_{k,:}.
\end{equation}

Here \(A_{:,k}\) is a column of length \(M\) and \(B_{k,:}\) a row of
length \(N\); their outer product has shape \texttt{{[}M,N{]}}. For
\(k=0\), the contribution is \([1,0,2]^{\mathsf T}[2,1]\), giving rows
\texttt{{[}2,1{]}}, \texttt{{[}0,0{]}}, \texttt{{[}4,2{]}}. Adding all
four outer products yields the same \(C\). Row-dot-column explains a
cell, repeated GEMV explains columns, and outer products explain reuse
across cells. None changes the mathematical result over real numbers.

Tracing all six cells matters. A test that checks only \(C_{00}\) can
miss a wrong output-row stride. A square-only test can let code confuse
\(K\) with \(N\) without crossing a bound. ENGINE-2 includes asymmetric
fixtures and a grid of rectangular shapes for exactly this reason.

\section{Build it: reference GEMM}\label{build-it-reference-gemm}

The reference algorithm mirrors the equation:

\begin{verbatim}
for i in 0..rows {
    for j in 0..columns {
        let mut sum = 0.0_f32;
        for k in 0..inner {
            sum += *left.get2(i, k)? * *right.get2(k, j)?;
        }
        output[i * columns + j] = sum;
    }
}
\end{verbatim}

Calling this loop \emph{reference} is more precise than dismissing it as
merely naive. It prioritizes correspondence with the mathematical
specification and uses checked logical access. That makes it useful for
review, unusual strides, and differential testing. It is intentionally
not the performance endpoint.

The output allocation is checked before numerical work. A result shape
\texttt{{[}M,N{]}} whose element count overflows \texttt{usize} returns
\texttt{\seqsplit{OutputShapeOverflow}}. Valid zero cases remain
defined:

\begin{itemize}
\tightlist
\item
  \texttt{{[}M,0{]}\ ×\ {[}0,N{]}} returns \texttt{{[}M,N{]}} filled
  with zeros;
\item
  \texttt{{[}0,K{]}\ ×\ {[}K,N{]}} returns an empty \texttt{{[}0,N{]}}
  owner;
\item
  \texttt{{[}M,K{]}\ ×\ {[}K,0{]}} returns an empty \texttt{{[}M,0{]}}
  owner.
\end{itemize}

The implementation checks representability, but ordinary \texttt{Vec}
allocation can still fail under genuine memory exhaustion. Rust's
standard global allocator behavior is outside this small API's
recoverable error model; the distinction between arithmetic overflow and
resource exhaustion must remain explicit.

\begin{quote}
\textbf{BUILD IT} Complete
\href{https://github.com/hermonai/inside-the-llm-engine/blob/astra-visual-rewrite/labs/lab-24-gemm-by-hand.md}{Lab
24}. Check every cell and then use a transpose view to prove the
reference path follows strides.
\end{quote}

\section{How much arithmetic?}\label{how-much-arithmetic}

There are \(MN\) output cells and \(K\) multiply-accumulate
contributions per cell. Counting one multiplication and one addition as
two floating-point operations gives approximately

\begin{equation}
F_{\mathrm{GEMM}}\approx2MKN\quad[\mathrm{FLOPs}].
\end{equation}

If one counts the first assignment differently, the exact addition count
is \(MN(K-1)\) for nonempty \(K\), while the multiply count is \(MNK\).
The conventional \(2MKN\) expression is a useful performance model, not
a statement that every compiler emits exactly that many instructions.

For GEMV,

\begin{equation}
F_{\mathrm{GEMV}}\approx2MK\quad[\mathrm{FLOPs}].
\end{equation}

\textbf{FLOP} names an amount of floating-point work. \textbf{FLOP/s}
names a rate. An observed effective rate based on elapsed time \(t\)
seconds is

\begin{equation}
P_{\mathrm{effective}}
\approx\frac{2MKN}{t}\quad[\mathrm{FLOP/s}].
\end{equation}

The benchmark reports GFLOP/s by dividing that rate by \(10^9\). It does
not claim those operations correspond one-for-one to retired hardware
instructions.

\section{Why equal FLOPs do not mean equal
time}\label{why-equal-flops-do-not-mean-equal-time}

Two loop nests can compute the same \(2MKN\) model and take different
times. The processor must obtain operands before multiplying them and
retain or reload partial outputs before adding. Arithmetic units are
only one resource among caches, load/store units, memory channels,
address-generation machinery, and front-end control.

\Appref{tensors} measured a simpler clue on this same Apple M1. Traversing one
\texttt{{[}2048,2048{]}} allocation in row-major order had a 4,163,875
ns median; visiting the same values column-wise had a 13,692,583 ns
median. That experiment was not GEMM and cannot prove a matrix-kernel
ratio. It established a narrower fact: access order can matter even when
the values and arithmetic are held constant. \Appref{matmul} therefore
measures the matrix loops directly.

\section{Follow the operand}\label{follow-the-operand}

In canonical row-major storage, flat offsets are

\begin{equation}
\operatorname{offset}_A(i,k)=iK+k,
\end{equation}

\begin{equation}
\operatorname{offset}_B(k,j)=kN+j,
\end{equation}

and

\begin{equation}
\operatorname{offset}_C(i,j)=iN+j.
\end{equation}

For fixed \(i\) and \(j\), the reference \texttt{i,j,k} loop increments
\(k\). Consecutive \(A_{ik}\) values have adjacent offsets. Consecutive
\(B_{kj}\) values are \(N\) elements apart because the loop walks down
one logical column of row-major \(B\). The
\href{https://github.com/hermonai/inside-the-llm-engine/blob/astra-visual-rewrite/diagrams/linear/row-major-access.txt}{row-major
access diagram} contrasts those movements.

The arithmetic is symmetric in the pair of operands; their memory walks
are not. When \(N\) is large, asking the innermost loop to step by \(N\)
can underuse each fetched cache line. The machine may retrieve
neighboring \(B\) values that this particular inner loop does not
immediately consume.

\section{Spatial and temporal
locality}\label{spatial-and-temporal-locality}

\textbf{Spatial locality} means accessing nearby addresses close
together in time. A cache line normally transfers several adjacent
values, so a stride-1 walk can use more of the transferred line.
\textbf{Temporal locality} means reusing the same value or region before
it leaves a fast level of the memory hierarchy.

These are tendencies, not commands to hardware. Prefetchers, cache
mapping, compiler transformations, matrix dimensions, and competing work
affect the outcome. Still, they give us a productive question: can a
legal loop permutation make the intended reuse easier to realize?

\section{Working set and the memory
hierarchy}\label{working-set-and-the-memory-hierarchy}

\EngineFigure{ch06-hierarchy}{Conceptual memory-to-execution hierarchy with spatial and temporal locality distinguished.}{fig:ch06-hierarchy}

A processor does not generally load every scalar directly from main
memory on every use. Registers sit closest to arithmetic. Several cache
levels retain recently accessed lines. Main memory is larger and usually
more expensive to reach. Exact capacities, policies, and latencies are
architecture-specific, but the hierarchy creates a general objective:
arrange work so values are reused while they remain in a faster level.

The \textbf{working set} is the data actively needed during a region of
execution. For one output cell in \texttt{ijk}, the conceptual set
includes an A row, a B column, and one accumulator. The A row has a
contiguous walk, but the B column touches many separated cache lines
when N is large. Across adjacent output columns, the algorithm revisits
the same A row, yet whether it remains resident depends on the rest of
the active footprint.

In \texttt{ikj}, one \(A_{ik}\) scalar is loaded into a local variable
and reused across an output row. The active B row and C row are
contiguous. This does not ensure that an entire long row fits a
particular cache, but it exposes short-range spatial reuse to the
hardware and compiler.

Blocking shrinks that focus again. Instead of streaming complete rows
while other operands compete for capacity, it asks the machine to
revisit bounded A, B, and C regions. A useful tile lets multiple
arithmetic contributions occur per line fetched from a lower level. A
bad tile can be too large to retain, too small to amortize control, or
poorly shaped for the matrix.

Cache lines also explain why counting logical scalar loads in source is
not a traffic measurement. Loading one \texttt{f32} can bring
neighboring values. A later access may hit that line without another
main-memory transfer, or the line may have been evicted and fetched
again. The minimum-byte equations describe data that must logically
participate; only appropriate performance counters and a defined
hierarchy boundary could establish physical bytes.

This distinction prevents a common reasoning error. We may say
\texttt{ikj} presents contiguous B access and that blocking is designed
to increase temporal reuse. We may not say exactly how many cache misses
it removes on the recorded host, because the current benchmark records
elapsed time rather than cache events.

There are six permutations of three distinct loop indices:

\begin{verbatim}
ijk   ikj   jik   jki   kij   kji
\end{verbatim}

Not all can use a single scalar accumulator in the same way, and their
access patterns differ under row-major layout. ENGINE-2 studies
\texttt{ijk} and \texttt{ikj} rather than pretending to rank all six
universally.

\section{\texorpdfstring{The \texttt{ikj}
rewrite}{The ikj rewrite}}\label{the-ikj-rewrite}

The update equation can be written as

\begin{equation}
C_{ij}\mathrel{+}=A_{ik}B_{kj}.
\end{equation}

For a fixed \(i\) and \(k\), let \(j\) vary through a row segment:

\begin{verbatim}
for i in 0..m {
    for k in 0..inner {
        let left_value = left[i * inner + k];
        for j in 0..n {
            output[i * n + j] += left_value * right[k * n + j];
        }
    }
}
\end{verbatim}

Now the innermost loop walks adjacent \(B_{kj}\) and \(C_{ij}\)
positions. One loaded \(A_{ik}\) scalar feeds an entire output-row
segment. The canonical
\href{https://github.com/hermonai/inside-the-llm-engine/blob/astra-visual-rewrite/diagrams/linear/loop-order-ijk-vs-ikj.txt}{loop-order
diagram} and
\href{https://github.com/hermonai/inside-the-llm-engine/blob/astra-visual-rewrite/diagrams/linear/follow-the-reuse.txt}{reuse
diagram} show this change.

Initialization moved too. \texttt{ijk} can compute one local sum and
assign a cell. \texttt{ikj} accumulates into many cells across K
iterations, so the output must be zero-initialized before the loop.
Forgetting that step produces dependence on old storage rather than
matrix multiplication.

\begin{quote}
\textbf{PERFORMANCE LAB}
\href{https://github.com/hermonai/inside-the-llm-engine/blob/astra-visual-rewrite/labs/lab-25-loop-order.md}{Lab
25} begins with physical offset sequences, not a stopwatch. Correctness
is the admission ticket to timing.
\end{quote}

\section{Performance lab: loop
order}\label{performance-lab-loop-order}

\EngineFigure{ch06-loops}{IJK strides down B while IKJ sweeps adjacent B and C elements, preserving all contributions.}{fig:ch06-loops}

For the fixture, canonical strides are \texttt{A:{[}4,1{]}},
\texttt{B:{[}2,1{]}}, \texttt{C:{[}2,1{]}}. At \texttt{i=0,j=0}, IJK
reads B offsets \texttt{0,2,4,6}; IKJ at \texttt{i=0,k=0} reads offsets
\texttt{0,1} and updates C offsets \texttt{0,1} using the same A scalar.
The
\href{https://github.com/hermonai/inside-the-llm-engine/blob/astra-visual-rewrite/figures/generated/ch06-loops.html}{loop-order
animation} advances \(k\); its static plate shows all four states. The
parity gate verifies that the full IJK and IKJ traces contain the
identical set of 24 \texttt{(i,j,k)} contributions. This proves
traversal coverage, not a cache-hit count or a speedup.

The release harness compares direct scalar row-major \texttt{ijk} and
\texttt{ikj} loops. Both allocate and zero their result, use
deterministic \texttt{f32} inputs, and pass a per-element correctness
gate. On the recorded Apple M1 run at code commit \texttt{03e08a87}
(full pin in the record), median results were:

{\def\LTcaptype{none} % do not increment counter
\begin{longtable}[]{@{}rrrrr@{}}
\toprule\noalign{}
Square size & Repetitions & \texttt{ijk} & \texttt{ikj} & Observed
ratio \\
\midrule\noalign{}
\endhead
\bottomrule\noalign{}
\endlastfoot
64 & 15 & 185,375 ns & 54,791 ns & 3.38× \\
128 & 9 & 1,645,000 ns & 286,375 ns & 5.74× \\
256 & 5 & 15,638,708 ns & 1,709,000 ns & 9.15× \\
\end{longtable}
}

The exact environment, command, raw checksums, compiler, and limitations
live in the
\href{https://github.com/hermonai/inside-the-llm-engine/blob/astra-visual-rewrite/research/benchmarks/chapter-06-loop-order.md}{loop-order
benchmark record}. For these tested shapes, \texttt{ikj} had the lower
median. The widening ratio is consistent with its contiguous inner
walks; it is not a portable law or a substitute for hardware-counter
evidence.

\section{What the benchmark does---and does
not---measure}\label{what-the-benchmark-doesand-does-notmeasure}

The \Appref{matmul} harness uses \texttt{\seqsplit{std::time::Instant}} in a
Cargo release build. Inputs are deterministically generated before
timing. Every candidate is executed once for correctness before samples
are accepted. Timed output is consumed through a checksum and
\texttt{black\_box}, reducing the chance that an optimizing compiler
discards the work. Samples are sorted and the median is reported.

Median is a defensible summary for this bounded, exploratory
microbenchmark, but it does not describe tail latency or serving
variability. Repetition counts appear beside every result. Small shapes
use more repetitions because each sample is short; larger shapes use
fewer to keep the exercise practical. The process is warm, but cache
contents are not flushed or controlled between samples. Candidate order
is stable rather than randomized. Frequency, temperature,
operating-system activity, and compiler-generated instructions are not
independently measured.

Allocation and zero-initialization are included. That matches the
current public API, which returns a fresh \texttt{OwnedTensor}, and
prevents a prepared-output control from receiving a hidden advantage. It
also means the record cannot attribute every nanosecond to arithmetic or
cache traversal. A future reusable output API would require a separate
experiment with equivalent initialization policy.

The harness's direct \texttt{ijk} and \texttt{ikj} functions use
canonical flat slices after their deterministic construction. Timing
\texttt{\seqsplit{matmul\_reference}} itself would also include a
checked logical view lookup for every operand access, confounding the
loop-order question with repeated boundary validation. Correctness still
uses the public contracts; the loop-order experiment isolates the two
flat scalar traversals. The blocked experiments time the actual public
blocked kernel, including its metadata views and layout checks.

Effective GFLOP/s divides the conventional work estimate by elapsed
time. It is useful for comparing shapes inside this record, but it is
not measured peak hardware throughput. The ideal FLOP/byte column
divides work by compulsory logical payload, not hardware counters.
Publishing both without these labels would create precision without
accuracy.

Finally, the benchmark is not an inference benchmark. It has no
tokenizer, request queue, model file, KV cache, sampling, streaming,
concurrency, or token latency. Its purpose is narrower: make one kernel
hypothesis reproducible. End- to-end claims require end-to-end workloads
and cannot be obtained by multiplying independent microbenchmark ratios.

\section{Memory movement}\label{memory-movement}

A simplistic implementation might reload \(A_{ik}\) and \(B_{kj}\) for
every mathematical contribution, then repeatedly read and write
\(C_{ij}\). Caches can remove some transfers to lower memory levels, but
only if the active data stays resident and the access order exposes
reuse.

The smallest compulsory payload model for canonical \texttt{f32} GEMM
reads each input once and writes each output once:

\begin{equation}
Q_{\min}=4(MK+KN+MN)\quad[\mathrm{bytes}].
\end{equation}

This is a lower-bound model, not measured traffic. Write allocation,
eviction, cache-line granularity, repeated loads, and metadata can
increase physical movement. The
\href{https://github.com/hermonai/inside-the-llm-engine/blob/astra-visual-rewrite/diagrams/linear/follow-the-byte.txt}{byte-flow
diagram} keeps the ownership and transfer story separate from the
formula.

\section{Arithmetic intensity}\label{arithmetic-intensity}

\textbf{Arithmetic intensity} is floating-point work divided by bytes
moved at a chosen memory boundary:

\begin{equation}
I=\frac{F}{Q}\quad[\mathrm{FLOP/byte}].
\end{equation}

Using the compulsory-payload model for GEMM gives

\begin{equation}
I_{\mathrm{ideal}}
\approx\frac{2MKN}{4(MK+KN+MN)}\quad[\mathrm{FLOP/byte}].
\end{equation}

The word \emph{ideal} is essential. Without hardware counters, ENGINE-2
does not know actual cache or DRAM traffic.

For GEMV with \(A:[M,K]\), \(x:[K]\), and \(y:[M]\), the analogous model
is

\begin{equation}
I_{\mathrm{GEMV,ideal}}
\approx\frac{2MK}{4(MK+K+M)}\quad[\mathrm{FLOP/byte}].
\end{equation}

For large \(M\) and \(K\), the \(MK\) weight term dominates, so
intensity approaches approximately \(0.5\) FLOP/byte. One weight
contributes to one output vector. In GEMM, a weight can contribute
across \(N\) columns, so potential reuse and ideal intensity grow with
\(N\).

\subsection{Future connection: prefill and
decode}\label{future-connection-prefill-and-decode}

Later chapters will distinguish two inference workload phases.
\textbf{Prefill} processes many prompt positions and can present
matrix-like collections of activations. \textbf{Decode} advances one or
a few new positions per active sequence and can present vector-like or
skinny-matrix work. That suggests a useful first intuition:

Many activation columns can reuse a weight across multiple positions; one column offers less such reuse. This is a shape-dependent opportunity, not a hardware-independent performance law.

This is explicitly a conceptual preview, not a complete workload model.
Attention shapes, batching, KV state, quantization, providers, and
scheduling will refine it. A production runtime may also batch decode
work from several requests into a matrix. \Appref{matmul} establishes only the
numerical reason that the number of activation columns can change
weight-reuse opportunity.

\section{The Roofline mental model}\label{the-roofline-mental-model}

The Roofline model relates attainable performance to peak compute,
memory bandwidth, and arithmetic intensity:

\begin{equation}
P\le\min\left(P_{\mathrm{peak}},\ B_{\mathrm{memory}}I\right).
\end{equation}

At low intensity, the sloped bandwidth term can be the tighter ceiling.
At high intensity, the horizontal compute ceiling can dominate. The
canonical
\href{https://github.com/hermonai/inside-the-llm-engine/blob/astra-visual-rewrite/diagrams/linear/roofline-concept.txt}{Roofline
concept diagram} is a mental model, not a calibrated roofline for this
laptop. We have not measured peak bandwidth, peak arithmetic throughput,
or per-level traffic.

This preview answers the chapter's central question: multiplication may
be cheap relative to repeatedly delivering its operands. An optimization
can improve performance without changing FLOPs if it improves realized
data reuse. Conversely, a theoretically higher-intensity workload can
remain slow when a poor scalar kernel fails to exploit that opportunity.

\section{Why blocking exists}\label{why-blocking-exists}

Loop interchange improves the immediate access direction, but whole rows
or matrices can still exceed a fast cache. \textbf{Blocking}, also
called \textbf{tiling}, partitions M, K, and N into bounded ranges. The
kernel works on an \(A\) tile, a \(B\) tile, and an accumulating \(C\)
tile before moving on.

The
\href{https://github.com/hermonai/inside-the-llm-engine/blob/astra-visual-rewrite/diagrams/linear/cache-reuse-and-tiling.txt}{tiling
diagram} shows the intended active set. For block extents
\((B_M,B_K,B_N)\), one rough \texttt{f32} payload is

\begin{equation}
Q_{\mathrm{tile}}=4(B_MB_K+B_KB_N+B_MB_N)\quad[\mathrm{bytes}].
\end{equation}

That formula does not select a universally correct tile. Cache capacity
is not the only constraint: associativity, cache lines, register
pressure, loop overhead, compiler decisions, and matrix shape all
matter. ENGINE-2 makes the three extents explicit and offers
\texttt{{[}32,32,32{]}} as a teaching default, not a machine truth.

\section{Build it: blocked scalar
GEMM}\label{build-it-blocked-scalar-gemm}

The blocked path uses outer tile loops \texttt{ii,kk,jj} and inner
scalar loops \texttt{i,k,j}. This is the actual loop body from
\texttt{linear.rs}, after validation and zero-initialized output
allocation:

\begin{verbatim}
for ii in (0..rows).step_by(block.m) {
    let i_end = ii.saturating_add(block.m).min(rows);
    for kk in (0..inner).step_by(block.k) {
        let k_end = kk.saturating_add(block.k).min(inner);
        for jj in (0..columns).step_by(block.n) {
            let j_end = jj.saturating_add(block.n).min(columns);
            for i in ii..i_end {
                let output_row = i * columns;
                let left_row = i * inner;
                for k in kk..k_end {
                    let left_value = left_data[left_row + k];
                    let right_row = k * columns;
                    for j in jj..j_end {
                        output[output_row + j] += left_value * right_data[right_row + j];
                    }
                }
            }
        }
    }
}
\end{verbatim}

The real implementation uses saturating addition before \texttt{min} so
even endpoint arithmetic cannot wrap. Each block extent must be
positive; otherwise \texttt{step\_by(0)} would be invalid and the
intended partition would be meaningless.

Tail tiles are mandatory. For \texttt{{[}5,7{]}\ ×\ {[}7,3{]}} with
blocks \texttt{{[}4,4,2{]}}, all three axes have a partial final tile. A
benchmark that tests only dimensions evenly divisible by the block size
can report a fast but incomplete kernel.

\begin{quote}
\textbf{BUILD IT}
\href{https://github.com/hermonai/inside-the-llm-engine/blob/astra-visual-rewrite/labs/lab-26-blocked-gemm.md}{Lab
26} traces those three tails and compares every result cell with the
reference path.
\end{quote}

\section{Separate checked boundary from hot
loop}\label{separate-checked-boundary-from-hot-loop}

\EngineFigure{ch06-tiles}{Rectangular tile offsets and clipped edge regions in a five-by-seven times seven-by-three product.}{fig:ch06-tiles}

The tail fixture deliberately uses a different shape from the small
numerical spine: \texttt{{[}5,7{]}\ ×\ {[}7,3{]}} with blocks
\texttt{{[}4,4,2{]}}. The last tile covers \texttt{i=4..5},
\texttt{k=4..7}, \texttt{j=2..3}, all half-open. Its three scalar
contributions are as real as those in a full tile. There are eight tile
combinations; clipping bounds avoids padding requirements and dropping
remainders. The test compares every output against the general
reference. Tile diagrams show flat offsets here, not invented matrix
values.

The blocked kernel does not call checked multidimensional indexing for
each multiply. It validates rank, inner dimensions, output size, block
size, and layout once.
\texttt{\seqsplit{TensorView::as\_contiguous\_slice}} then returns the
exact logical range only when strides equal canonical row-major strides.
Inside the loop, flat indices operate on safe Rust slices.

This is not the removal of safety. It is the movement of repeated
validation to a boundary whose proof applies to the whole loop. The
crate retains \texttt{\seqsplit{\#![forbid(unsafe\_code)]}}; slice
indexing remains bounds checked by Rust, and shape construction has
already proved representable storage extents.

Reference and blocked policies are intentionally different:

\Needspace{14\baselineskip}

{\def\LTcaptype{none} % do not increment counter
\begin{longtable}[]{@{}
  >{\raggedright\arraybackslash}p{(\linewidth - 4\tabcolsep) * \real{0.3333}}
  >{\raggedright\arraybackslash}p{(\linewidth - 4\tabcolsep) * \real{0.3333}}
  >{\raggedright\arraybackslash}p{(\linewidth - 4\tabcolsep) * \real{0.3333}}@{}}
\toprule\noalign{}
\begin{minipage}[b]{\linewidth}\raggedright
Property
\end{minipage} & \begin{minipage}[b]{\linewidth}\raggedright
Reference
\end{minipage} & \begin{minipage}[b]{\linewidth}\raggedright
Blocked
\end{minipage} \\
\midrule\noalign{}
\endhead
\bottomrule\noalign{}
\endlastfoot
Accepted rank & matrices only & matrices only \\
Valid layouts & any checked read-only strides & strict canonical
row-major \\
Inner order & \texttt{i,j,k} & tiled \texttt{i,k,j} \\
Input mutation & never & never \\
Output & new canonical owner & new canonical owner \\
Hidden materialization & never & never \\
\end{longtable}
}

The
\href{https://github.com/hermonai/inside-the-llm-engine/blob/astra-visual-rewrite/diagrams/linear/reference-vs-blocked-kernel.txt}{kernel-contract
diagram} captures the boundary. If a caller wants to multiply a
transpose view with the blocked path, it must call
\texttt{to\_contiguous()} explicitly. The copy is then visible in code,
cost accounting, and ownership.

\begin{quote}
\textbf{ENGINEERING FAILURE --- THE INVISIBLE COPY} An optimized API
receives a strided transpose, silently packs it, runs a fast kernel, and
reports only kernel time. The result is correct but the system paid
allocation and copy costs that its call site cannot see. ENGINE-2
returns \texttt{\seqsplit{UnsupportedLayout}} instead.
\end{quote}

\section{Reference versus
optimized}\label{reference-versus-optimized}

An optimized substitution is valid only if it preserves the selected
contract. It need not preserve internal iteration order. For each
accepted input it must produce the same output shape and numerically
equivalent values, leave inputs unchanged, return independent owned
storage, handle empty dimensions, and return documented errors before
invalid execution.

Keeping the reference path makes those claims testable. It also leaves a
fallback for valid strided views and a small implementation that
reviewers can compare directly with the equation. Production libraries
keep analogous fallbacks and differential oracles because optimized code
has more edge paths, not fewer.

The
\href{https://github.com/hermonai/inside-the-llm-engine/blob/astra-visual-rewrite/diagrams/linear/optimization-ladder.txt}{optimization
ladder} places blocking in context. Packing, register blocking, SIMD
microkernels, threading, NUMA placement, accelerators, and fusion remain
later rungs. ENGINE-2 stops at scalar cache blocking so the first
optimization remains inspectable. The
\href{https://github.com/hermonai/inside-the-llm-engine/blob/astra-visual-rewrite/diagrams/linear/reference-candidate-gates.txt}{reference/candidate
gate} states the contract, equivalence, and performance-evidence
sequence required before a candidate can replace the oracle path.

\section{Prove it: independent
oracle}\label{prove-it-independent-oracle}

\texttt{\seqsplit{code/reference/python/chapter06\_matmul\_oracle.py}}
imports no Rust or mini-engine code. It implements dot, GEMV, and GEMM
with Python lists and uses \texttt{struct.pack}/\texttt{unpack} to round
each multiplication and addition to binary32. Its fixtures include:

\begin{itemize}
\tightlist
\item
  the exact dot result 12 for \texttt{{[}1,2,3{]}\ ·\ {[}4,-5,6{]}};
\item
  the asymmetric GEMV result \texttt{{[}1.5,6{]}};
\item
  the complete hand GEMM result \texttt{\seqsplit{[[58,64],[139,154]]}};
\item
  fractional values and negative terms;
\item
  a logical transpose product;
\item
  the empty-dot identity.
\end{itemize}

An independent oracle can share a mistaken specification, but it is less
likely to share the same indexing defect than a second call through the
Rust kernel. Hand expansion supplies another line of evidence.

\section{Prove it: equivalence gate}\label{prove-it-equivalence-gate}

ENGINE-2 uses

\begin{equation}
|a-r|\le \epsilon_{abs}+\epsilon_{rel}|r|,
\end{equation}

with \(\epsilon_{abs}=10^{-5}\) and \(\epsilon_{rel}=10^{-5}\) for
chapter fixtures. Absolute tolerance protects values near zero; relative
tolerance scales with result magnitude. The tolerance is narrow enough
for these bounded fixtures and is not declared universal for arbitrary
reduction lengths or value distributions.

The deterministic property-style test spans every \(M,K,N\) from 0
through 6 and three distinct block shapes. It includes empty products,
tiny matrices, exact tiles, and tails without introducing a random-test
dependency. Additional tests cover strided reference inputs, zero-stride
views, invalid ranks, mismatched inner dimensions, non-canonical blocked
operands, zero block dimensions, output overflow, and independent output
ownership.

\section{Floating-point order is
observable}\label{floating-point-order-is-observable}

Real-number addition is associative:

\begin{equation}
(a+b)+c=a+(b+c).
\end{equation}

Finite floating-point addition need not be. Rounding after an
intermediate sum can discard information that a different grouping
preserves. A simple reduction order is therefore part of a
reproducibility story even when it is not part of the abstract linear
algebra.

An executable binary32 example uses \(a=2^{24}\), \(b=1\),
\(c=-2^{24}\):

\begin{equation}
\operatorname{fl}(\operatorname{fl}(a+b)+c)=0,
\qquad
\operatorname{fl}(a+\operatorname{fl}(b+c))=1.
\end{equation}

At that positive magnitude, the first grouping loses the unit before
cancellation; the second preserves it. The new test evaluates both
groupings as \texttt{f32}, with black-box inputs. This is rounding, not
an indexing defect.

Fusing also changes the rounding points. Set \(u=1+2^{-13}\) and
\(v=1-2^{-13}\). Both are exactly representable binary32 inputs. Their
exact product is \(1-2^{-26}\), which rounds to 1 when stored as F32:

\begin{equation}
\operatorname{fl}(\operatorname{fl}(uv)-1)=0,
\qquad \operatorname{fma}(u,v,-1)=-2^{-26}.
\end{equation}

The test uses explicit \texttt{mul\_add} only to demonstrate this
distinction. The teaching kernels do not call it. A vector reduction may
create several partial sums and combine them later, adding another
reason to state numerical policy separately from the real-number
equation.

The reference kernel visits K in ascending order for each cell. The
blocked kernel visits K tiles in ascending order and K within each tile
in ascending order. Its interleaving across different output cells
changes, but the sequence of contributions to an individual cell remains
ordered. That makes the current paths especially close. The API still
specifies tolerance rather than bitwise equivalence because compiler
contraction, target instructions, and future legal kernel
transformations can affect rounding.

Using \texttt{f64} accumulation would reduce error for many \texttt{f32}
inputs, but it would also define a different kernel cost and numerical
contract. ENGINE-2 chooses \texttt{f32} input, accumulator, and output
because it keeps the scalar teaching path aligned with its stored dtype.
The choice is explicit, not asserted optimal. Later low-precision or
quantized kernels will have to state storage, product, and accumulator
types separately.

NaN and infinity policy also deserves precision. Model construction
rejects non-finite parameters, and \texttt{Logits} rejects non-finite
results. The generic linear kernels themselves perform ordinary
\texttt{f32} arithmetic and do not scan all inputs for finiteness. That
keeps them mathematical primitives rather than silently embedding the
model's parameter policy. A caller that needs a finite- input contract
must validate at its own boundary.

Determinism is correspondingly bounded. For the same build, target,
inputs, and selected scalar path, iteration order is fixed. ENGINE-2
does not promise bit-identical output across all compilers and
architectures. It promises shape and ownership semantics plus numerical
equivalence under the documented test tolerance.

\begin{quote}
\textbf{PROVE IT} Run
\href{https://github.com/hermonai/inside-the-llm-engine/blob/astra-visual-rewrite/labs/lab-28-kernel-equivalence.md}{Lab
28}. Record maximum error; never widen a tolerance until the discrepancy
has an explanation.
\end{quote}

Floating-point tolerance cannot excuse a structural bug. A wrong offset
may occasionally produce a close number. That is why shape checks, exact
integer fixtures, full-output comparisons, tails, and error variants
surround the numerical metric.

\section{Performance lab: blocking}\label{performance-lab-blocking}

The blocked benchmark sweeps cubic tile sizes 8, 16, 24, 32, 48, and 64
at shape \texttt{192³}. On the recorded Apple M1 run, tile 64 had the
lowest median among those candidates, while tile 8 was slower than the
direct \texttt{ijk} control. The teaching default of 32 was not the
measured winner.

A second experiment held the 32 tile fixed and varied square size:

{\def\LTcaptype{none} % do not increment counter
\begin{longtable}[]{@{}
  >{\raggedleft\arraybackslash}p{(\linewidth - 6\tabcolsep) * \real{0.2500}}
  >{\raggedleft\arraybackslash}p{(\linewidth - 6\tabcolsep) * \real{0.2500}}
  >{\raggedleft\arraybackslash}p{(\linewidth - 6\tabcolsep) * \real{0.2500}}
  >{\raggedleft\arraybackslash}p{(\linewidth - 6\tabcolsep) * \real{0.2500}}@{}}
\toprule\noalign{}
\begin{minipage}[b]{\linewidth}\raggedleft
Size
\end{minipage} & \begin{minipage}[b]{\linewidth}\raggedleft
Direct \texttt{ijk}
\end{minipage} & \begin{minipage}[b]{\linewidth}\raggedleft
Blocked 32
\end{minipage} & \begin{minipage}[b]{\linewidth}\raggedleft
Speedup: direct time / blocked time
\end{minipage} \\
\midrule\noalign{}
\endhead
\bottomrule\noalign{}
\endlastfoot
8 & 333 ns & 1,000 ns & 0.33× \\
16 & 2,459 ns & 2,834 ns & 0.87× \\
32 & 19,375 ns & 12,500 ns & 1.55× \\
64 & 178,333 ns & 93,125 ns & 1.91× \\
128 & 1,668,500 ns & 755,625 ns & 2.21× \\
\end{longtable}
}

\EngineFigure{ch06-crossover}{Historical direct-time over blocked-time ratios retain the two losses below parity.}{fig:ch06-crossover}

For sizes 8 and 16, the blocked call lost on this run. Extra
tile/control work is a plausible explanation, but no experiment isolated
that cause. Blocking began winning somewhere between the tested sizes 16
and 32; that bracket is workload- and machine-specific. See the
\href{https://github.com/hermonai/inside-the-llm-engine/blob/astra-visual-rewrite/research/benchmarks/chapter-06-blocked-matmul.md}{blocked-matmul
record} for raw results and full limitations.

This is what an honest optimization result looks like: it contains
losing cases. ``Blocked'' is not a synonym for ``fast.'' Optimization
has a workload.

\section{Performance lab: GEMV versus
GEMM}\label{performance-lab-gemv-versus-gemm}

The third experiment holds a \texttt{{[}512,512{]}} weight matrix
constant while changing the right-hand column count. One column uses
scalar GEMV; 8 and 64 columns use blocked GEMM.

{\def\LTcaptype{none} % do not increment counter
\begin{longtable}[]{@{}rlrrr@{}}
\toprule\noalign{}
Columns \(N\) & Kernel & Median & Effective GFLOP/s & Ideal FLOP/byte \\
\midrule\noalign{}
\endhead
\bottomrule\noalign{}
\endlastfoot
1 & GEMV & 191,416 ns & 2.739 & 0.498 \\
8 & blocked GEMM & 3,294,459 ns & 1.273 & 3.879 \\
64 & blocked GEMM & 6,217,167 ns & 5.397 & 25.600 \\
\end{longtable}
}

\EngineFigure{ch06-throughput}{Measured GEMV and blocked GEMM throughput beside separately modeled ideal intensity.}{fig:ch06-throughput}

The analytic reuse opportunity rises monotonically. The observed scalar
rate did not. The narrow \texttt{N=8} blocked workload achieved less
effective throughput than the GEMV case, while \texttt{N=64} exceeded
both. Short inner row segments, tile overhead, and cache behavior are
hypotheses for this untuned kernel, not measured causes. The
\href{https://github.com/hermonai/inside-the-llm-engine/blob/astra-visual-rewrite/research/benchmarks/chapter-06-gemv-vs-gemm.md}{GEMV/GEMM
record} contains the reproducer and caveats; the
\href{https://github.com/hermonai/inside-the-llm-engine/blob/astra-visual-rewrite/diagrams/linear/gemv-vs-gemm-reuse.txt}{reuse
comparison diagram} shows only opportunity, not guaranteed performance.

\begin{quote}
\textbf{PERFORMANCE LAB} In
\href{https://github.com/hermonai/inside-the-llm-engine/blob/astra-visual-rewrite/labs/lab-29-gemv-vs-gemm.md}{Lab
29}, keep total latency, effective throughput, ideal bytes, and measured
bytes conceptually separate.
\end{quote}

\section{ENGINE-2: replace the projection
loop}\label{engine-2-replace-the-projection-loop}

All three performance sections above retain the \textbf{2026-09-03
historical record}, pinned to \texttt{03e08a87} (full hash in each
linked record). The regenerated charts parse those records directly.
Running the examples during this edition's verification does not replace
the published medians. The historical benchmark helper accepts
absolute/relative tolerances \texttt{1e-4/1e-5}; the bounded unit
fixtures use \texttt{1e-5/1e-5}. They are deliberately reported as
separate gates. Neither permits NaN to pass by a comparison trick.

The ENGINE-1 model now creates a rank-1 borrowed view over its
request-owned hidden activation and calls

\begin{verbatim}
gemv_reference(&self.output_weight.view(), &hidden_view)
\end{verbatim}

The kernel returns \(W\mathbf{h}\) as a new owned vector. The model then
adds its bias explicitly and constructs finite \texttt{Logits}. The
\href{https://github.com/hermonai/inside-the-llm-engine/blob/astra-visual-rewrite/diagrams/linear/weight-orientation.txt}{weight-orientation
diagram} shows that each \texttt{{[}V,D{]}} row belongs to one candidate
token. No transpose is implied or materialized.

Embedding remains a row lookup:

\begin{equation}
\mathbf{h}=E[x].
\end{equation}

Selecting row \(x\) from \(E:[V,D]\) is not matrix multiplication.
Recasting it as a one-hot GEMV would add \(VD\) mostly useless
arithmetic and obscure the actual operation. Kernel reuse should clarify
the model, not force every operation through one abstraction.

The \Appref{smallest-model} fixture remains the numerical regression gate. Input
\texttt{like} still produces logits

\begin{equation}
[-0.7,\ 0.1,\ 0.4,\ 2.2],
\end{equation}

and the autoregressive sequence remains \texttt{Rust} followed by EOS.
This proves that extracting GEMV changed architecture without changing
observable model semantics.

The
\href{https://github.com/hermonai/inside-the-llm-engine/blob/astra-visual-rewrite/diagrams/linear/engine-2-kernel-stack.txt}{ENGINE-2
stack diagram} shows the new boundary. Tensors describe storage; the
\texttt{linear} module performs operators; the model composes an
embedding lookup, GEMV, and bias addition. Putting \texttt{matmul}
methods directly on \texttt{OwnedTensor} would blur those layers and
make policy---reference versus blocked, layout requirements, future
provider selection---look like an intrinsic property of storage.

\section{The public kernel
contract}\label{the-public-kernel-contract}

ENGINE-2 exposes four free functions:

\begin{verbatim}
dot_reference(left, right) -> Result<f32, KernelError>
gemv_reference(matrix, vector) -> Result<OwnedTensor, KernelError>
matmul_reference(left, right) -> Result<OwnedTensor, KernelError>
matmul_blocked(left, right, block) -> Result<OwnedTensor, KernelError>
\end{verbatim}

\texttt{KernelError} distinguishes wrong rank, vector length mismatch,
matrix inner dimension mismatch, unsupported optimized layout, invalid
block size, output shape overflow, and lower-level tensor failure. The
error variants preserve operation and operand context where that context
helps diagnosis.

The contract has no output aliasing. Both inputs are immutable views.
Every tensor output is a fresh canonical owner. Reference calls accept
any view whose non-negative strides and extent pass Tensor Substrate v1.
Blocked calls require strict canonical row-major strides, including axes
of length one. Neither path silently changes shape or layout.

\begin{quote}
\textbf{BREAK IT}
\href{https://github.com/hermonai/inside-the-llm-engine/blob/astra-visual-rewrite/labs/lab-27-break-the-kernel.md}{Lab
27} turns every contract clause into a typed failure. A valid transpose
is accepted by reference GEMM and rejected by blocked GEMM until the
caller explicitly materializes it.
\end{quote}

\section{Inside Hermon}\label{inside-hermon-2}

Hermon is the book's industrial reference, not the source of ENGINE-2's
API. At freshly inspected commit \texttt{2a3fd521} (full pin in the
source ledger), its safe Rust surface validates matrix/vector and
matrix/matrix dimensions before crossing a foreign-function boundary.
The linked runtime routes packed weights into a native tensor bridge.
That bridge constructs GGML tensors and an operation graph around
\texttt{ggml\_mul\_mat}; a tensor session owns execution state across
the unsafe boundary.

Status matters. Hermon's default execution remains its batched llama.cpp
path. Its Hermon-owned paged path is PREVIEW and gated rather than
silently selected. The \texttt{paged.rs} row-major \texttt{f32}
multiplication is test/oracle support, not proof that the production
runtime uses a Rust scalar GEMM for normal inference. The GGUF packed
path validates layout and shapes before native graph construction.

\begin{quote}
\textbf{INSIDE HERMON --- CURRENT/LIBRARY/PREVIEW} CURRENT Hermon checks
contracts and dispatches its default batched runtime. LIBRARY code in
pinned GGML performs the industrial matrix operation. PREVIEW
Hermon-owned paging remains explicitly gated. These layers must not be
blended into one performance claim.
\end{quote}

The comparison is useful because it exposes how much ENGINE-2 postpones:
packed weight formats, native graph lifetime, backend dispatch,
architecture specialization, threads, and accelerators. It also
validates the architectural lesson that layout and shape checks belong
before the hot native execution boundary.

The current default route is concrete: \texttt{dispatch.rs} selects the
batched runtime; \texttt{batched.rs} calls
\texttt{\seqsplit{Context::decode\_batch}}; \texttt{linked.rs} crosses
\texttt{\seqsplit{hermon\_llama\_decode\_batch}}; \texttt{shim.c}
invokes \texttt{llama\_decode}. The model graph then reaches GGML
operations and configured backends. This is not the same route as the
optional packed tensor bundle API.

For that separate API, \texttt{matvec\_f32} delegates to
\texttt{matmul\_f32} with one input row. The bundle call validates one
through eight weight names, positive input rows, compatible
rank/dimensions and shared width, checked element counts and native
integer conversions. A mutable \texttt{TensorSession} supplies exclusive
access to reusable context/threadpool state. The C++ bridge checks host
weight residency, copies the F32 input into a GGML tensor, creates one
\texttt{ggml\_mul\_mat} node per weight, computes the graph, and copies
results into Rust-owned output vectors. It is a \textbf{host CPU
bridge}, not evidence that this helper dispatches to the default model's
GPU backend. \texttt{paged.rs} uses the bundle interface on the gated
PREVIEW path; library availability alone does not make it the default.

\EngineFigure{ch06-source}{Default batched execution and optional host bridge converge on GGML semantics, with independent ENGINE-2 verification below.}{fig:ch06-source}

The
\href{https://github.com/hermonai/inside-the-llm-engine/blob/astra-visual-rewrite/research/astra/chapter06-regeneration.md}{dated
source ledger} records the exact files, symbols, commits, and status of
these findings. Local source inspection establishes available paths and
checks, not a claim that a particular GPU ran a particular request. This
chapter did not launch an industrial model or benchmark Hermon.

\section{Inside llama.cpp and GGML}\label{inside-llama.cpp-and-ggml}

Hermon's pinned llama.cpp submodule was inspected at \texttt{389ff61d}
(full pin in the source ledger). GGML's public \texttt{ggml\_mul\_mat}
contract uses its own dimension conventions and supports selected type
pairs. CPU dispatch selects type-specific vector-dot functions. The CPU
kernel examines strides and contiguity, partitions work into
tiles/chunks, and can select architecture-specific implementations;
other backends have separate Metal and CUDA \texttt{MUL\_MAT} paths.

The production lesson is not ``copy GGML into the teaching engine.'' It
is that high-performance multiplication is a family of shape-, type-,
layout-, and hardware-dependent kernels behind a validated contract.
Even names and matrix orientation must be translated carefully at a
boundary. ENGINE-2's
\texttt{{[}M,K{]}\ ×\ {[}K,N{]}\ -\textgreater{}\ {[}M,N{]}} convention
is deliberately explicit so readers do not infer semantics from a
foreign function name.

GGML also demonstrates the optimization rungs beyond this chapter:
type-driven dot kernels, packing assumptions, register-scale work,
parallel scheduling, and accelerator dispatch. Those mechanisms are
real, but introducing them before a scalar oracle would make failures
harder to localize.

\EngineFigure{ch06-production}{Teaching tensor contracts compared with GGML axis, byte-stride, and packed-type conventions.}{fig:ch06-production}

Make the orientation translation explicit. The teaching equation uses
\(W:[M,K]\), \(B:[K,N]\), and \(C:[M,N]\). The bundle bridge instead
receives row-batched \(X=B^{\mathsf T}:[N,K]\) and returns
\(Y=XW^{\mathsf T}=C^{\mathsf T}:[N,M]\). In GGML's dimension order the
weight has \texttt{ne={[}K,M{]}}, the input \texttt{ne={[}K,N{]}}, and
the result \texttt{ne={[}M,N{]}}:

\begin{equation}
Y_{ji}=\sum_{k=0}^{K-1} X_{jk}W_{ik}=C_{ij}.
\end{equation}

These are coordinate correspondences, not permission to pass the
teaching \(B\) buffer unchanged. Canonical teaching \(B\) interleaves
columns; canonical row-batched \(X\) stores each input vector together.
For the fixture, \(X\) has rows \texttt{{[}2,1,-1,0.5{]}} and
\texttt{{[}1,-1,2,0{]}}, and \(Y\) has rows \texttt{{[}-3,3,2.5{]}} and
\texttt{{[}9,-3,4{]}}. Preparing that layout is an explicit boundary
responsibility. An axis label may look reversed while the contraction is
correct; a copied buffer with the wrong order can have correct
dimensions and wrong values.

GGML's \texttt{ne{[}{]}} describes logical extents and \texttt{nb{[}{]}}
byte strides. For a packed quantized type, a storage block may hold
multiple logical weights plus scale metadata. Its byte size is not four
times its logical element count. The CPU type traits choose a compatible
vector-dot function and right-hand dot type; the implementation may
convert the right operand into work storage before computing. This
previews representation-dependent kernels, not a quantization format,
accuracy result, or implemented ENGINE-2 feature.

\section{From scalar tiles to CPU
microkernels}\label{from-scalar-tiles-to-cpu-microkernels}

Blocking the cache working set is not the last reuse opportunity. A CPU
microkernel can hold a small output rectangle in registers while
stepping through K. One loaded A value can be broadcast across vector
lanes; a vector of adjacent B values contributes to several output
columns. This is the IKJ reuse idea at register scale. A dot-oriented
kernel instead builds several partial sums and eventually reduces them
horizontally. SIMD does not mean one instruction completes the entire
matrix; the finite register set, vector width, dependency chains, and
edge handling still constrain the schedule.

The pinned \texttt{\seqsplit{ggml-cpu/vec.cpp}} makes the distinction
observable in source: \texttt{\seqsplit{ggml\_vec\_dot\_f32}} has a
\texttt{GGML\_SIMD} path, vector loads, FMA macro operations, multiple
accumulator vectors and a final reduction. \texttt{ggml-cpu.c} selects
\texttt{vec\_dot} and \texttt{vec\_dot\_type} from CPU type traits and
partitions matrix work into chunks. This is inspected implementation
evidence, not a promise that every build enables the same instruction
set or that every shape takes that path. The book's scalar Rust loops
remain a separate reference.

Packing can arrange input panels for such a microkernel, replacing
awkward strides with its preferred address pattern. It adds a copy,
storage, and lifetime cost. Persistent model weights can amortize
preparation across calls, while a small transient matrix may never repay
it. Threading adds another independent dimension: output regions can be
partitioned, but placement, shared-cache use, and synchronization
influence performance. BLIS documents loop-level thread partitioning
around a microkernel and separates build-time thread support from
runtime selection; the presence of a threaded library is not proof that
a call used several cores.
\href{https://github.com/flame/blis/blob/master/docs/Multithreading.md}{BLIS
threading documentation}

\section{GPU execution is a different schedule, not a different
equation}\label{gpu-execution-is-a-different-schedule-not-a-different-equation}

\EngineFigure{ch06-hardware}{Conceptual CPU vector work and GPU cooperative tile work, explicitly outside ENGINE-2 implementation scope.}{fig:ch06-hardware}

A GPU assigns pieces of an output tile to cooperating groups of threads.
A common tiled design loads operand regions from device/global memory,
stages reusable data in shared or threadgroup storage, and keeps partial
outputs in registers. Thread indexing must make collective memory
accesses efficient; barriers order cooperative reuse where required.
NVIDIA's worked matrix multiplication examples distinguish improved
global access from eliminating redundant loads. These are distinct
benefits, not evidence collected by this chapter.
\href{https://docs.nvidia.com/cuda/cuda-c-best-practices-guide/index.html\#shared-memory-in-matrix-multiplication-c-ab}{CUDA
Best Practices, shared-memory matrix multiplication}

CUTLASS explains a hierarchy of thread-block, warp, and thread tiles,
with register accumulators inside the larger tile. Specialized matrix
multiply-accumulate instructions can consume matrix fragments, including
mixed-precision combinations, rather than one scalar product per source
statement. Storage, multiplication, and accumulation types must then be
specified independently. This is a conceptual preview, not a claim that
all GPU kernels use one staging scheme or numerical type.
\href{https://developer.nvidia.com/blog/cutlass-linear-algebra-cuda/}{NVIDIA's
CUTLASS explanation}

At the pinned GGML revision, CUDA has distinct floating-point/quantized
matrix-vector and matrix-matrix branches and a cuBLAS route. Metal's
\texttt{\seqsplit{ggml\_metal\_op\_mul\_mat}} examines dimensions,
operand types and eligibility; its small-batch branches differ from its
matrix-matrix paths. The literal thresholds in that source belong to
that implementation, not to ENGINE-2's Apple M1 crossover experiment. A
CUDA warp and Metal SIMD group are not a portable synonym for the
chapter's four drawn lanes.

The scheduling problem includes launch cost, data residency, transfers
when needed, occupancy, register pressure, tails, and synchronization.
More nominal arithmetic capacity cannot rescue every tiny or skinny
workload. Conversely, large regular work can amortize costs and exploit
reuse unavailable to one isolated vector. Neither statement supplies a
speedup number. Measure the actual shape, representation, backend,
precision and end-to-end boundary.

\section{Why this is still not production
GEMM}\label{why-this-is-still-not-production-gemm}

ENGINE-2 is dependency-free, scalar, single-threaded, CPU-only, and safe
Rust. It does not pack panels, use SIMD intrinsics, specialize register
microkernels, call BLAS, use Apple's Accelerate framework, spawn Rayon
workers, dispatch by CPU feature, fuse bias, or offload to a GPU. It
also allocates a new output on each public call.

Production GEMM commonly layers:

\begin{enumerate}
\def\labelenumi{\arabic{enumi}.}
\tightlist
\item
  semantic validation and transpose/layout interpretation;
\item
  packing into kernel-friendly panels;
\item
  cache blocking;
\item
  register blocking and vector microkernels;
\item
  thread partitioning and NUMA policy;
\item
  architecture and accelerator dispatch;
\item
  fusion with adjacent operations where semantics permit.
\end{enumerate}

The full
\href{https://github.com/hermonai/inside-the-llm-engine/blob/astra-visual-rewrite/diagrams/linear/optimization-ladder.txt}{optimization
ladder diagram} marks exactly where the milestone stops. Previewing
later rungs is useful; implementing them now would violate the chapter's
controlled experiment.

Packing deserves particular caution. A packed layout can make a
microkernel fast but costs time and memory to produce. Immutable model
weights may amortize packing across many calls; one-off activations may
not. That accounting belongs with the future kernel that actually packs.
The same principle applies to thread launch and GPU transfer: overhead
is part of the workload, not an asterisk after the timing.

\section{Engineering failures}\label{engineering-failures}

\subsection{Inner dimensions do not
agree}\label{inner-dimensions-do-not-agree}

Code receives \texttt{{[}M,K₁{]}\ ×\ {[}K₂,N{]}} and loops to one
operand's bound. It either indexes outside the other view or computes a
truncated fiction. ENGINE-2 returns
\texttt{\seqsplit{InnerDimensionMismatch}} before allocation.

\subsection{The right offset uses the wrong
dimension}\label{the-right-offset-uses-the-wrong-dimension}

Writing \texttt{right{[}k\ *\ K\ +\ j{]}} instead of
\texttt{right{[}k\ *\ N\ +\ j{]}} can pass square tests. Asymmetric
\texttt{{[}2,3{]}\ ×\ {[}3,2{]}} and property grids expose it.

\subsection{Tile tails disappear}\label{tile-tails-disappear}

Loops assume dimensions divide the tile exactly. Boundary rows or
columns stay zero, or an unchecked implementation crosses storage. The
\texttt{{[}5,7{]}\ ×\ {[}7,3{]}} fixture with \texttt{{[}4,4,2{]}}
blocks creates tails on every axis.

\subsection{Faster but wrong}\label{faster-but-wrong}

A benchmark times an optimized path without comparing output. Dead work,
missing tiles, or stale accumulation can look exceptionally fast. The
harness checks outputs first and consumes checksums so the result
remains observable.

\subsection{Allocation is counted on one side
only}\label{allocation-is-counted-on-one-side-only}

One candidate reuses a prepared output while another allocates and
zeros. The comparison attributes memory-management policy to loop order.
The \Appref{matmul} harness includes output allocation for both direct
variants and documents the blocked API's allocation.

\subsection{Debug performance is
reported}\label{debug-performance-is-reported}

Rust debug builds prioritize diagnostics and overflow checks rather than
optimized execution. Chapter performance records use
\texttt{cargo\ run\ -\/-release} and identify compiler and flags.

\subsection{FLOPs become
instructions}\label{flops-become-instructions}

The conventional \texttt{2MKN} count is described as retired operations,
even though FMA, vector width, compiler transformations, and loop
control make that false. The book calls it an effective work model.

\subsection{Tolerance hides a structural
defect}\label{tolerance-hides-a-structural-defect}

A generous threshold makes wrong indexing ``equivalent.'' ENGINE-2
combines a tight bounded tolerance with exact hand fixtures, shape
assertions, typed failures, and full-vector comparison.

\subsection{Benchmark order becomes cache
policy}\label{benchmark-order-becomes-cache-policy}

One candidate always runs after another and inherits warmer inputs. The
small harness uses a warm process but does not fully randomize or
control cache state; the records disclose this limitation. Serving
claims would require a stronger protocol and distributions, not just
medians from this microbenchmark.

\section{Common misconceptions, checked against this
chapter}\label{common-misconceptions-checked-against-this-chapter}

\begin{itemize}
\tightlist
\item
  \textbf{``GEMV is trivial.''} Its equation is compact; the physical
  GEMV trace still requires correct orientation, checked offsets,
  ownership and a reduction.
\item
  \textbf{``GEMM is just more GEMV.''} Repeated columns prove
  mathematical equivalence; the outer-product view exposes a different
  opportunity to reuse each weight.
\item
  \textbf{``More FLOPs means slower.''} The historical N sweep separates
  total work, latency and attained throughput. None alone predicts the
  other two.
\item
  \textbf{``Asymptotic complexity selects the fastest loop.''} IJK and
  IKJ have the same contribution count, yet their recorded medians
  differ. Constants and address order remain observable below the
  asymptotic notation.
\item
  \textbf{``Blocking always helps.''} The crossover record includes
  losses at sizes 8 and 16, and the tile sweep includes a losing tile 8.
\item
  \textbf{``The default block size is optimal.''} Tile 64, not the
  default 32, won the recorded 192-cube sweep. One sweep cannot
  establish a portable default.
\item
  \textbf{``SIMD changes the equation.''} The same contributions can
  occupy vector lanes; scheduling and F32 grouping change, not the
  intended contraction.
\item
  \textbf{``Matrix multiplication is one instruction.''} The
  source-to-kernel map distinguishes a graph node from loads,
  arithmetic, stores and dispatch.
\item
  \textbf{``A GPU multiplies everything at once.''} The hardware plate
  shows finite tiles, registers and cooperative groups; launch and data
  movement still cost.
\item
  \textbf{``Equivalent kernels must be bit-identical.''} The executable
  reassociation and FMA examples show why numerical equivalence needs a
  stated policy.
\item
  \textbf{``Transpose always copies.''} \Appref{tensors}'s metadata view
  remains valid for the reference; blocked layout rejection makes any
  later copy explicit.
\item
  \textbf{``Quantization changes matrix shape.''} GGML's logical extents
  survive a representation change; bytes and approximation error require
  separate rules.
\item
  \textbf{``llama.cpp defines multiplication.''} The equation precedes
  the library. The source comparison translates its axes and dispatch,
  not the mathematics.
\end{itemize}

\section{Labs and exercises}\label{labs-and-exercises}

The chapter's executable sequence is:

\begin{itemize}
\tightlist
\item
  \href{https://github.com/hermonai/inside-the-llm-engine/blob/astra-visual-rewrite/labs/lab-22-dot-product.md}{Lab
  22 --- Dot Product}, CHECK;
\item
  \href{https://github.com/hermonai/inside-the-llm-engine/blob/astra-visual-rewrite/labs/lab-23-gemv-by-hand.md}{Lab
  23 --- GEMV by Hand}, CHECK;
\item
  \href{https://github.com/hermonai/inside-the-llm-engine/blob/astra-visual-rewrite/labs/lab-24-gemm-by-hand.md}{Lab
  24 --- GEMM by Hand}, BUILD;
\item
  \href{https://github.com/hermonai/inside-the-llm-engine/blob/astra-visual-rewrite/labs/lab-25-loop-order.md}{Lab
  25 --- Loop Order}, BUILD;
\item
  \href{https://github.com/hermonai/inside-the-llm-engine/blob/astra-visual-rewrite/labs/lab-26-blocked-gemm.md}{Lab
  26 --- Blocked GEMM}, BUILD;
\item
  \href{https://github.com/hermonai/inside-the-llm-engine/blob/astra-visual-rewrite/labs/lab-27-break-the-kernel.md}{Lab
  27 --- Break the Kernel}, BREAK;
\item
  \href{https://github.com/hermonai/inside-the-llm-engine/blob/astra-visual-rewrite/labs/lab-28-kernel-equivalence.md}{Lab
  28 --- Kernel Equivalence}, CHECK;
\item
  \href{https://github.com/hermonai/inside-the-llm-engine/blob/astra-visual-rewrite/labs/lab-29-gemv-vs-gemm.md}{Lab
  29 --- GEMV Versus GEMM}, EXTEND.
\end{itemize}

Further exercises:

\begin{enumerate}
\def\labelenumi{\arabic{enumi}.}
\tightlist
\item
  Enumerate the physical \(B\) offsets for all six loop orders on
  \texttt{{[}2,3{]}\ ×\ {[}3,4{]}}. Classify only the innermost access;
  do not predict a universal winner.
\item
  Derive the exact multiplication and addition counts for \(K=0\),
  \(K=1\), and \(K>1\). Compare them with the conventional \(2MKN\)
  model.
\item
  For a hypothetical 32 KiB cache, solve for cubic \texttt{f32} tile
  sizes whose ideal three-tile payload fits. List at least three reasons
  the calculation does not guarantee residency or speed.
\item
  Add a deterministic cancellation-heavy fixture. Compare ascending and
  descending K reductions and explain the observed error.
\item
  Design an output-buffer API that permits reuse. State the aliasing,
  initialization, shape, and failure rules before writing code.
\item
  Explain when explicitly materializing a transpose once could be
  profitable across many products. Include the copy in a break-even
  model.
\item
  Extend the benchmark with rectangles where \(M\ll N\) and \(N\ll M\).
  Preserve negative results and avoid tuning a global default from one
  host.
\item
  Read the Netlib SGEMM interface and list the features ENGINE-2
  intentionally omits. Explain why omission makes this chapter's
  contract easier to prove.
\end{enumerate}

\section{What we still have not
built}\label{what-we-still-have-not-built-1}

At the \Appref{matmul} curriculum boundary, ENGINE-2 is a linear algebra
kernel layer, not a Transformer. This milestone adds no normalization,
learned Transformer block, attention, Q/K/V projections, causal mask,
RoPE, KV cache, GGUF loader, quantized storage, SIMD, BLAS integration,
threading, GPU provider, autograd, or distributed execution.

It also does not choose kernels dynamically. \texttt{gemv\_reference}
and \texttt{matmul\_blocked} are explicit calls with explicit contracts.
A future provider may use shape, dtype, layout, hardware, and measured
crossover information to select an implementation, but silently
inventing that policy here would make the first optimization cycle
harder to see.

\section{Summary}\label{summary-5}

Matrix multiplication is a family of reductions. Dot product reduces two
length-\(K\) vectors to one scalar. GEMV applies one row dot product per
output element. GEMM applies one row/column dot product per output cell:

\begin{equation}
C_{ij}=\sum_{k=0}^{K-1}A_{ik}B_{kj}.
\end{equation}

The conventional work models are \(2MK\) FLOPs for GEMV and \(2MKN\)
FLOPs for GEMM. Those counts do not determine elapsed time. Row-major
offset progression, spatial locality, temporal reuse, working-set size,
and kernel overhead all matter. Arithmetic intensity connects work to
modeled byte movement; the Roofline model previews why low-reuse
operations can be bandwidth constrained.

ENGINE-2 now provides transparent dot, GEMV, and GEMM reference kernels
plus a scalar blocked GEMM. The reference path supports valid strided
views. The blocked path requires canonical storage and never hides
materialization. Both return explicit owned outputs, use checked shape
boundaries, preserve zero-size semantics, and report typed failures.

The optimized path did not replace its oracle. Deterministic property
tests, tail fixtures, the independent Python implementation, numerical
tolerances, and the preserved ENGINE-1 logits form its equivalence gate.
Release measurements then showed both wins and losses: loop order
mattered on the recorded host, blocking hurt tiny matrices, tile choice
mattered, and reuse opportunity did not automatically become throughput.

\begin{quote}
\textbf{FIRST PRINCIPLE} Count the arithmetic, trace the bytes, prove
the substitution, and only then believe the stopwatch.
\end{quote}

\section{Preview of Chapter~\ref{app:embeddings-norm}}\label{chapter-7-preview}

The completed boundary can be read as a compact execution map:

Explicit tensor storage supports reference linear kernels; verified linear kernels then support embedding and normalization composition. Each layer retains the preceding oracle.

ENGINE-2 means transparent reference semantics over explicit tensor
storage, an independently verified blocked candidate, and bounded
measurements of locality-sensitive work. It is more than four function
names and less than a production backend.

We can now store tensors honestly and apply checked linear
transformations. The next missing primitive is normalization. \Appref{embeddings-norm}
will build RMSNorm from first principles: squares, mean square, epsilon,
reciprocal square root, learned scale, precision, and its place around
model blocks.

That chapter will reuse ENGINE-2's discipline---reference semantics
before optimization---but it does not begin here. Matrix multiplication
remains the only new numerical operator family in this milestone.

In the book's progression, \Appref{tensors} establishes where values live;
\Appref{matmul} shows how computation walks them; \Appref{embeddings-norm} uses parameters
and activations in embedding and normalization; \Appref{qkv} will give
learned projections Q/K/V roles. The existing \Appref{embeddings-norm} code remains a
regression, not newly regenerated work in this edition.

\section{References}\label{references-3}

\begin{itemize}
\tightlist
\item
  Netlib, \href{https://www.netlib.org/blas/}{Basic Linear Algebra
  Subprograms}, including the Level 2 GEMV and Level 3 GEMM families.
\item
Netlib LAPACK documentation,
  \href{https://netlib.org/lapack/explore-html/d7/dda/group__gemv_ga0d35d880b663ad18204bb23bd186e380.html}{SGEMV
  interface and reference source}.
\item
Netlib LAPACK documentation,
  \href{https://www.netlib.org/lapack/explore-html/d4/de2/sgemm_8f_source.html}{SGEMM
  reference source}.
\item
  Samuel Williams, Andrew Waterman, and David Patterson,
  \href{https://www2.eecs.berkeley.edu/Pubs/TechRpts/2008/EECS-2008-134.pdf}{Roofline:
  An Insightful Visual Performance Model for Multicore Architectures}.
\item
  Intel,
  \href{https://www.intel.com/content/www/us/en/docs/advisor/cookbook/2023-0/optimize-memory-access-patterns.html}{Optimize
  Memory Access Patterns}, for loop interchange and cache blocking
  examples.
\item
  Oracle,
  \href{https://docs.oracle.com/cd/E37069_01/html/E39019/gnabm.html}{When
  \texttt{(-x)\ *\ (-y)} Is Not \texttt{x\ *\ y}}, for floating-point
  reassociation constraints.
\item
  Project research note,
  \href{https://github.com/hermonai/inside-the-llm-engine/blob/astra-visual-rewrite/research/part-02/chapter-06-matrix-multiplication-the-engine-room.md}{Chapter~\ref{app:matmul} --- Matrix Multiplication: The Engine Room}, with commit-pinned
  Hermon and GGML paths.
\item
  Project benchmark records:
  \href{https://github.com/hermonai/inside-the-llm-engine/blob/astra-visual-rewrite/research/benchmarks/chapter-06-loop-order.md}{loop
  order},
  \href{https://github.com/hermonai/inside-the-llm-engine/blob/astra-visual-rewrite/research/benchmarks/chapter-06-blocked-matmul.md}{blocked
  multiplication}, and
  \href{https://github.com/hermonai/inside-the-llm-engine/blob/astra-visual-rewrite/research/benchmarks/chapter-06-gemv-vs-gemm.md}{GEMV
  versus GEMM}.
\end{itemize}


\clearpage
\section{Worked synthesis: explain the work and the traffic}
\subsection{A rectangular tail is not padding}
\textbf{Problem.} Multiply $A:[5,7]$ and $B:[7,3]$ with tile widths
$(T_M,T_K,T_N)=(4,4,2)$. What are the clipped extents of the final tile?

\textbf{Solution.} The last starts are $(i_0,k_0,j_0)=(4,4,2)$.
The extents are $(1,3,1)$. Iterating the nominal $(4,4,2)$ sizes at this
boundary accesses nonexistent rows and columns. Clip ends to $(5,7,3)$.
There are $5\cdot7\cdot3=105$ multiply-accumulate contributions overall;
the usual convention counts approximately 210 FLOPs.

\subsection{An ideal lower bound on traffic}
\textbf{Problem.} For F32 $A:[M,K]$, $B:[K,N]$, and freshly written $C:[M,N]$,
give the ideal intensity if each input is read once and each output written once.

\textbf{Solution.}
\begin{equation}
I_{\mathrm{ideal}}=\frac{2MKN}{4(MK+KN+MN)}
\quad[\mathrm{FLOP/byte}].
\end{equation}
This is an optimistic traffic model, not measured DRAM traffic. It excludes
repeated reads, write allocation, metadata and copies. A timing cannot establish
this byte count; use counters or a clearly identified analytical model.

\subsection{A failed optimization is a result}
\textbf{Problem.} The blocked kernel is slower on a small matrix but agrees
with the oracle. What should the experiment report?

\textbf{Solution rubric.} Retain both times, workload, compiler, host, warmup,
repetition count and statistic. Report the slowdown and investigate overhead
relative to useful work. Do not discard the small shape or claim blocking
always wins. Correctness admits a candidate to measurement; it does not
guarantee a speedup.

